% =============================================================
%  Section 01: 研究背景与动机
% =============================================================
\section{研究背景与动机}

\begin{frame}{研究背景：RAG 为什么仍值得做}
  \begin{columns}[T,onlytextwidth]
    \column{0.58\textwidth}
    \begin{block}{大语言模型的三类现实短板}
      \begin{itemize}
        \item \hlb{知识静态}：训练截止后难以持续更新。
        \item \hlb{事实不稳}：长尾实体与多跳问题易出现幻觉。
        \item \hlb{不可追溯}：答案往往缺少显式证据链。
      \end{itemize}
    \end{block}
    \begin{exampleblock}{RAG 的核心价值}
      通过\hl{检索外部证据}约束生成，兼顾
      \hlb{正确性}、\hlb{可更新性}与\hlb{可解释性}。
    \end{exampleblock}

    \column{0.42\textwidth}
    \begin{alertblock}{真实应用中的新矛盾}
      \begin{itemize}
        \item 知识源不止文本，还有\hl{表格}与\hl{图谱}
        \item 单轮检索难覆盖\hl{第二跳证据}
        \item 证据不足时，系统应学会\hl{拒答}
      \end{itemize}
    \end{alertblock}
  \end{columns}
\end{frame}

\begin{frame}{本文关注的三个核心问题}
  \begin{columns}[T,onlytextwidth]
    \column{0.33\textwidth}
    \begin{block}{\numbox{1}\ 异构接入}
      文本、表格、图谱能否进入\hl{统一检索入口}，并保留可追溯来源？
    \end{block}

    \column{0.33\textwidth}
    \begin{block}{\numbox{2}\ 按需补检索}
      当首轮证据不足时，能否用\hl{低成本方式}提升多跳覆盖率？
    \end{block}

    \column{0.33\textwidth}
    \begin{block}{\numbox{3}\ 安全拒答}
      当证据冲突或缺失时，系统如何在\hl{可用性}与\hl{安全性}间平衡？
    \end{block}
  \end{columns}

  \vspace{0.5em}
  \begin{exampleblock}{研究立场}
    本文不是单纯追求更高分数，而是把 RAG 视为\hl{正确率、成本与安全性共同约束下的系统问题}。
  \end{exampleblock}
\end{frame}
